\documentclass[../../../../SoftwareRequirementsSpecification.tex]{subfiles}
\begin{document}
% Richiede le macro definite in src/macros/useCaseMacros.tex
% (incluse nel preambolo di SoftwareRequirementsSpecification.tex)

\ucstart
\begin{tabularx}{\textwidth}{|l|X|}
  \hline
  \ucid{PT-08} \\ \hline
  \ucname{Visualizzazione Atleti} \\ \hline
  \ucactor{PT} \\ \hline
  \ucdescription{Il PT visualizza l'elenco dei propri
    atleti e ne consulta i dettagli.} \\ \hline
  \ucprecond{Il PT deve essere autenticato.} \\ \hline
  \ucpostcond{Il PT visualizza i dettagli dell'atleta
    selezionato.} \\ \hline
  \ucmainflow{
    \item Il PT atterra sulla pagina dell'elenco dei propri atleti.
    \item Il sistema mostra l'elenco degli atleti seguiti dal PT, in
      ordine alfabetico per cognome, riportando per ciascuno nome,
      cognome ed email.
    \item Il PT preme su un atleta dell'elenco.
    \item Il sistema mostra la pagina di dettaglio dell'atleta, che
      riporta:
      \begin{itemize}
        \item i dati anagrafici (nome, cognome, età, email);
        \item i dati fisici (altezza e peso);
        \item le palestre in cui l'atleta si allena;
        \item l'elenco delle schede assegnate, con nome, periodo di
          validità e stato;
        \item l'elenco delle sessioni di allenamento e delle richieste
          di sessione che riguardano l'atleta, con data, ora e stato.
      \end{itemize}
    \item Il PT può premere il tasto "Assegna nuova Scheda" per
      assegnare una nuova scheda all'atleta.
    \item Il flusso continua dal passo 2 del caso d'uso PT-05.
  } \\ \hline
  \ucaltflow{
    \textbf{2a. Nessun atleta presente} \newline
    - Il sistema mostra un messaggio ("Nessun atleta registrato"). \newline
    - Il flusso termina.
  } \\ \hline
  \ucnote{
    \textbf{Nota 1:} tutte le informazioni della pagina di dettaglio
    sono in sola lettura. Lo statement attribuisce al PT la sola
    visualizzazione dei dettagli dei propri atleti, mentre
    l'inserimento e la modifica dei dati anagrafici spettano
    all'Atleta (caso d'uso A-05).\newline
    \textbf{Nota 2:} gli elenchi di schede, sessioni e richieste
    possono essere vuoti; in tal caso il sistema mostra la
    corrispondente sezione vuota, senza che il flusso cambi.\newline
    \textbf{Nota 3:} l'elenco delle schede assegnate è mostrato nel
    dettaglio anche perché consente al PT di verificare, prima di
    assegnarne una nuova tramite PT-05, che gli intervalli temporali
    delle schede dell'atleta non si sovrappongano.\newline
    \textbf{Nota 4:} il PT vede esclusivamente i propri atleti;
    l'elenco corrisponde agli atleti associati al PT autenticato.
  } \\ \hline
\end{tabularx}
\end{document}
